\section{Methods \& Models}
\label{sec:methods}

\paragraph{Architecture.} A two-layer MLP with GELU and dropout maps frozen
CLIP ViT-B/16 features ($\mathbb{R}^{512}$) to a $d$-dimensional embedding.
A geometry-specific head produces the final embedding: identity for the
Euclidean case, exponential map at the origin for the Poincar\'e ball,
\(\exp_0^c(u) = \tanh(\sqrt{c}\,\|u\|) \, u / (\sqrt{c}\,\|u\|)\). A
prototype classifier with one learnable point per style yields logits as
$-d(z, p_k)$, where $d(\cdot,\cdot)$ is Euclidean or hyperbolic accordingly.

\paragraph{Geometry, distance, and optimization.} The Poincar\'e distance
with curvature $c > 0$ is
\[
  d^c(x, y) = \frac{1}{\sqrt{c}}\,\operatorname{arcosh}\!\left(
    1 + \frac{2c\,\|x-y\|^2}{(1-c\|x\|^2)(1-c\|y\|^2)}
  \right).
\]
Standard linear layers do not preserve the ball, so the prototype
classifier is the natural choice: logits are distances, with no need for a
parametric ``softmax head'' inside the manifold. The MLP head is optimized
with Adam; the hyperbolic prototypes are optimized with Riemannian Adam
\citep{kochurov2020geoopt}, which projects each update back onto the ball.

\paragraph{Hierarchy-aware loss.} The MS3 baseline trains with cross-entropy
only and uses the hand-built tree only at evaluation time. For MS4 we add a
regularizer that pushes the prototype-prototype distance matrix toward the
tree distance matrix:
\[
  \mathcal{L} = \mathcal{L}_{CE} + \lambda \cdot \frac{1}{|\mathcal{P}|}
    \sum_{(i,j) \in \mathcal{P}} \!\left(
      \frac{d(p_i, p_j)}{\bar{d}} - \frac{T_{ij}}{\bar{T}}
    \right)^{\!\!2}.
\]
The mean-normalization (each pair-distance vector is divided by its mean
before the squared difference) makes the term scale-invariant: the geometry
chooses its own absolute embedding scale rather than fighting the
classification objective for magnitude. $\lambda = 0$ recovers the MS3
baseline exactly. This regularizer is computed once per gradient step over
the $\binom{27}{2} = 351$ off-diagonal style pairs --- a negligible cost
relative to the classification forward/backward.

\paragraph{Ground-truth trees.} The default tree is a hand-built lineage
hierarchy following standard art-history references (Renaissance branches
into Baroque branches into Romanticism; see Appendix~\ref{app:tree}). We
additionally evaluate two variants: a \emph{chronological} two-level era
grouping (Renaissance / Baroque-Rococo / 19th century / Early 20th /
Modern / Non-Western) and a \emph{flat} tree (every leaf is a direct child
of the root, used as a null since flat tree distances are constant
off-diagonal and cannot distinguish models).

\paragraph{Training details.} All runs use the supplied 70/30 train/val
split ($\approx$57k / 24k images), batch size 4096, Adam with $\eta=10^{-3}$
for the head and $10^{-2}$ for the prototypes, weight decay $10^{-4}$,
dropout 0.1, gradient clipping at norm 1.0, 30 epochs (training is fast
because the CLIP features are pre-cached as float16 tensors). Each headline
configuration is run with three seeds. The full grid (Phases 1--3, see
Section~\ref{sec:results}) totals $\approx 170$ runs and completed in under
one hour of wall time on a single Apple M-series GPU.

\paragraph{Evaluation suite.} We report top-1, top-5, and balanced
accuracy; prototype-tree and class-center-tree Spearman correlation;
multiplicative tree distortion (mean and worst-case); dendrogram-cluster
F1 against the default tree; sibling and cousin recall@$k$ for
$k \in \{5, 10\}$; and Fr\'echet-mean nearest-prototype accuracy on Cubism
as an interpolation-stability probe. Implementation lives in
\texttt{scripts/eval.py}.
